\section{Experiments}
\label{sec:experiments}

\subsection{Setup}
\label{sec:setup}

\paragraph{Models and encoding.} We use Qwen3-4B, 8B and 14B~\citep{qwen3_2025}.
Each model is encoded with 131{,}072 calibration tokens of
DCLM-edu~\citep{dclmedu}, drawn once per size. It is then trained and
assembled as in \S\ref{sec:model}.

\paragraph{Quality.} The benchmark is MMLU~\citep{mmlu2021}, 5-shot, on all
14{,}042 test questions, and every question counts the same (micro-average).
The answer is read from the logits of the four answer letters. Each file is
scored on its dense reconstruction, run through an ordinary forward pass. The
kernel computes with the same decoded weights. The per-row checks of
\S\ref{sec:kernel} and the 256-token comparison below tie the two. Every comparison is paired on the same questions. Each
gap gets a 95\,\% interval from a bootstrap that resamples questions within
each subject (10{,}000 draws). An exact McNemar test uses the questions that
only one of the two models gets right. This test set cannot separate two
scores less than about 0.4 points apart.

\paragraph{Speed and memory.} Decode speed is measured on one NVIDIA L40S at
batch 1, with greedy sampling. Speed is the tokens generated divided by the
wall-clock time, prompt included. Each speed is the median of five timed
rounds after a warm-up. Our files run in our engine at the served settings and generate 256
tokens. FP16 and AWQ run in vLLM 0.26.0~\citep{vllm2023}, AWQ
with the Marlin kernel, and generate 128 tokens. IQ2\_XXS runs in llama.cpp's
\texttt{llama-bench} on 128 tokens from an empty prompt, as the mean of five
repetitions. Each engine is only compared with itself, and we never divide one
engine's speed by another's. Memory counts only weight bytes: our engine's GPU
buffers for projections and tables, and the other formats' weight files. It
leaves out the KV cache and the activations.

\paragraph{Baselines.} FP16 is Qwen's published checkpoint. AWQ is Qwen's
official w4g128 checkpoint at each size (4-bit weights in groups of 128). Its
MMLU is scored in our harness after converting it to f16, and its speed is
timed in vLLM on the packed checkpoint. IQ2\_XXS is built with \texttt{llama-quantize} and an
importance matrix, which tells the quantizer which weights matter most. That
matrix comes from the same calibration shard as our earlier GPTQ runs.
IQ2\_XXS is scored with the same prompts through \texttt{llama-server}. We
cite the MMLU of QTIP, QuIP\# and the original LLVQ from \citet{llvq2026}
and do not run them.

\subsection{Main results}
\label{sec:main}

\begin{table}[t]
\centering\small\setlength{\tabcolsep}{5pt}
\begin{tabular}{llrrrr}
\toprule
 & format and engine & b/param & MMLU & decode tok/s & weights GB\\
\midrule
\multirow{8}{*}{4B}
 & FP16, vLLM & 16.00 & 70.14 & 83.1 & 8.04\\
 & f16 dense path, our engine$^{a}$ & 16.00 & & 43.0 & 8.04\\
 & AWQ w4g128, vLLM & 5.30 & 68.14 & 200.5 & 2.67\\
 & IQ2\_XXS, llama.cpp & 2.48 & 39.78 & 312.9 & 1.25\\
 & \textbf{\Tetra{}, our engine} & \textbf{2.73} & \textbf{63.37} & \textbf{113.8} & \textbf{1.38}\\
 & \textcolor{okGray}{LLVQ, cited$^{b}$} & \textcolor{okGray}{2 b/w$^{c}$} & \textcolor{okGray}{62.8} & & \\
 & \textcolor{okGray}{QTIP, cited$^{b}$} & \textcolor{okGray}{2 b/w$^{c}$} & \textcolor{okGray}{59.5} & & \\
 & \textcolor{okGray}{QuIP\#, cited$^{b}$} & \textcolor{okGray}{2 b/w$^{c}$} & \textcolor{okGray}{52.9} & & \\
\midrule
\multirow{4}{*}{8B}
 & FP16, vLLM & 16.00 & 75.05 & 46.3 & 16.38\\
 & f16 dense path, our engine$^{a}$ & 16.00 & & 26.4 & 16.38\\
 & AWQ w4g128, vLLM & 5.96 & 73.79 & 123.2 & 6.10\\
 & \textbf{\Tetra{}, our engine} & \textbf{2.70} & \textbf{69.58} & \textbf{95.0} & \textbf{2.76}\\
\midrule
\multirow{4}{*}{14B}
 & FP16, vLLM & 16.00 & 78.88 & 25.8 & 29.54\\
 & f16 dense path, our engine$^{a}$ & 16.00 & & 16.9 & 29.54\\
 & AWQ w4g128, vLLM & 5.40 & 78.12 & 77.8 & 9.98\\
 & \textbf{\Tetra{}, our engine} & \textbf{2.73} & \textbf{75.66} & \textbf{57.2} & \textbf{5.04}\\
\bottomrule
\end{tabular}
\caption{Quality, speed and memory at three sizes, on one L40S. MMLU uses all
14{,}042 test questions. The speeds come from different engines, and we never
divide one engine's speed by another's. $^{a}$Our engine's dense path runs the sealed file
decoded to f16 as an ordinary model, in the same process as the kernel. Its
speed is that of any f16 model of this size in our engine, which copies the
output head at every token. This row has no MMLU of its own.
$^{b}$Fine-tuned rows of Table~6 of \citet{llvq2026}, measured in their harness
with their MMLU protocol. $^{c}$Bits per weight, on the projections only. The
source gives no rate for the whole model. Data: paper2-results.csv.}
\label{tab:main}
\end{table}

\begin{figure}[t]
\centering
\includegraphics{fig_scale.pdf}
\caption{Against model size. Left: paired MMLU gaps (FP16 minus our file, AWQ
minus our file, FP16 minus AWQ), with 95\,\% intervals. Middle: bits per
parameter over the whole model, our sealed files against the official AWQ
checkpoints. Right: decode speed at batch 1, each format in its own engine.
Data: paper2-gaps.csv, paper2-chain.csv, paper2-results.csv.}
\label{fig:scale}
\end{figure}

\paragraph{Quality.} The sealed files score 63.37, 69.58 and 75.66 at 4B, 8B
and 14B (Table~\ref{tab:main}). On the same questions, they sit below FP16 by
6.77 points $[6.05, 7.50]$, 5.48 $[4.83, 6.10]$ and 3.22 $[2.69, 3.75]$. They
sit below AWQ by 4.76 $[4.02, 5.49]$, 4.21 $[3.54, 4.86]$ and 2.46
$[1.92, 3.00]$. Both gaps shrink as the model grows
(Figure~\ref{fig:scale}, left). AWQ itself is 2.00, 1.27 and 0.75
points below FP16.

\paragraph{Memory.} The sealed files use 2.73, 2.70 and 2.73 bits per
parameter, against 5.30, 5.96 and 5.40 for AWQ. That is 45 to 52\,\% of AWQ's
memory (Figure~\ref{fig:scale}, middle). AWQ keeps the embedding and the
output head in f16. At 8B these two tables are 15.2\,\% of the parameters, so
AWQ needs the most bits per parameter there.

\paragraph{Against other 2-bit formats at 4B.} IQ2\_XXS scores 39.78 at 2.48
bits per parameter, 23.6 points below our 4B file for 0.25 fewer bits per
parameter. The cited rows come from another harness and count bits on the
projections only. So they show roughly where our 63.37 stands, without ranking
it. It is close to the 62.8 of the original LLVQ with fine-tuning, and above
the 59.5 of QTIP.

\paragraph{Speed.} In our engine, the sealed files decode 113.8, 95.0 and 57.2
tokens per second at 4B, 8B and 14B (Figure~\ref{fig:scale}, right). Their
weights take 1.38, 2.76 and 5.04~GB on the card. Our dense path runs their
dense reconstruction in the same processes, at 43.0, 26.4 and 16.9 tokens per
second. A large part of this gap comes from the 4-bit output head. With the
embedding and head in f16 on both sides, the sealed files decode 55.0, 38.9
and 27.3 tokens per second. This is 1.27, 1.47 and 1.61
times our dense path. In vLLM, AWQ decodes 200.5, 123.2 and 77.8 tokens per
second, and FP16 83.1, 46.3 and 25.8. We do not divide these by our numbers,
because vLLM runs FP16 faster than our dense path does. A ratio across the two
engines would mix the format with the engine.

\paragraph{The served tokens.} At 4B and 8B, each sealed file run through the
kernel and its dense reconstruction, in the same process, generate the same
256 greedy tokens.
At 14B they first differ at token 78. In a greedy decode, one changed token
changes every later one, so only the first difference carries information. Our
preregistration, written before measuring, counts a divergence after token 32
as a near tie between two logits, not a defect. A second check feeds 203
prompt tokens in one call and in 203 single-token calls. Both give the same
next token at all three sizes.

\subsection{Where the gain comes from}
\label{sec:chain}

\begin{table}[t]
\centering\small\setlength{\tabcolsep}{5pt}
\begin{tabular}{llrrrl}
\toprule
 & step & b/param & MMLU & gain & 95\,\% interval, McNemar\\
\midrule
\multirow{3}{*}{4B}
 & \Tetra{} $+$ int4 \texttt{v\_proj} & 2.7475 & 57.95 & & \\
 & $+$ trained row scales & 2.7475 & 61.11 & $+3.15$ & $[+2.56, +3.74]$, $p=1.8\cdot10^{-25}$\\
 & $+$ int4 $o$, down; 4-bit tables & 2.7320 & 63.37 & $+2.26$ & $[+1.66, +2.87]$, $p=8.2\cdot10^{-13}$\\
\midrule
\multirow{3}{*}{8B}
 & \Tetra{} $+$ int4 \texttt{v\_proj} & 3.0683 & 64.87 & & \\
 & $+$ trained row scales & 3.0683 & 68.16 & $+3.29$ & $[+2.78, +3.80]$, $p=7.2\cdot10^{-38}$\\
 & $+$ int4 down; 4-bit tables & 2.6953 & 69.58 & $+1.42$ & $[+0.93, +1.92]$, $p=3.2\cdot10^{-8}$\\
\midrule
\multirow{3}{*}{14B}
 & \Tetra{} $+$ int4 \texttt{v\_proj} & 2.7371 & 72.53 & & \\
 & $+$ trained row scales & 2.7371 & 74.20 & $+1.67$ & $[+1.27, +2.08]$, $p=5.0\cdot10^{-16}$\\
 & $+$ int4 $o$, down; 4-bit tables & 2.7305 & 75.66 & $+1.46$ & $[+0.99, +1.94]$, $p=1.4\cdot10^{-9}$\\
\bottomrule
\end{tabular}
\caption{How each file is built, one step per row. Each gain is paired against
the row above on the 14{,}042 questions. Bits per parameter count the file as
served, with 8-bit embedding tables in the first two rows of each size. Those
rows were scored on MMLU with f16 tables. So the last gain also includes the
cost of going to 4-bit tables. Data: paper2-chain.csv.}
\label{tab:chain}
\end{table}

Table~\ref{tab:chain} scores each step against the previous one. Training the
row scales gains 3.15, 3.29 and 1.67 points without adding a byte. Adding int4
projections and 4-bit tables gains 2.26, 1.42 and 1.46 points more, while the
files get smaller. This test set separates all six
gains from zero, with every $p$ below $10^{-7}$.

\paragraph{Which projections get the int4 bytes.} At 8B, our first candidate
for the sealed file stored all of \texttt{o\_proj}, plus \texttt{down\_proj}
of layers 2 to 33, in int4. It scored 70.08 but needed 3.06 bits per parameter, because the
8B model pays for two tables, the embedding and the output head. To reach 2.7,
we gave int4 a smaller budget and tried two ways to spend it. The first puts
all of \texttt{o\_proj}, plus \texttt{down\_proj} of layers 15 to 20, in int4,
and scores 68.81 at 2.7047 bits per parameter. The second leaves
\texttt{o\_proj} as it is, puts \texttt{down\_proj} of layers 10 to 26 in
int4, and scores 69.58 at 2.6953.
It beats the first by 0.77 points $[0.26, 1.28]$, $p=0.004$, with fewer bytes.
It is the 8B file of Table~\ref{tab:main}. It sits 0.50 points below the
3.06-bit file, and this test set cannot separate the two ($[-0.01, 1.01]$,
$p=0.063$). We had predicted the opposite order, from an earlier exploration
at 4B on a smaller MMLU sample. The 4B and 14B files still store
\texttt{o\_proj} in int4, and we have not tested \texttt{down\_proj} alone
there.

\paragraph{A second training does not help.} A second training on the sealed
4B file also learns the 73 RMSNorm weight vectors. It lowers the KL loss from
0.262 to about 0.20 and moves MMLU by $+0.02$ points $[-0.31, +0.36]$ ($p=0.93$). We keep the file
without it.
